\chapter{Table Formats and the Iceberg Spec}
\label{ch:iceberg-spec}

\begin{chapterintro}
An Iceberg table is not a directory of Parquet files with a manifest bolted on.
It is a tree: a catalog pointer to a metadata file, to a manifest list, to
manifests, to data and delete files --- every level immutable, every commit a
single atomic pointer swap. Reads consult this tree before touching any data.
This chapter walks the spec (format version 2) against a real table's metadata
tables so you can draw the tree from memory.
\end{chapterintro}

\section{From Hive tables to table formats}
\tbd

\section{The Iceberg spec, format version 2}
\tbd

\section{Snapshots and the metadata tree}
\tbd

\section{Position deletes versus equality deletes}
\tbd

\section{Iceberg, Delta Lake, and Hudi}
\tbd

\section{Summary}
\tbd

\exercisehead
\begin{exercises}
\item Commit three appends to a table, then inspect \texttt{metadata\_log\_entries},
      \texttt{snapshots}, and \texttt{manifests}. Explain how many metadata files
      and manifests exist and why.
\end{exercises}
